\documentclass[../main.tex]{subfiles}
\pagenumbering{roman}
\begin{document}
\begin{center}
    \large{\color{primary}\textbf{LỜI MỞ ĐẦU}}\\
\end{center}
\vspace{1cm}
Trong bối cảnh thương mại điện tử ngày càng phát triển mạnh mẽ, việc xây dựng các hệ thống phần mềm có khả năng mở rộng, dễ bảo trì và đáp ứng nhanh trước sự thay đổi của nghiệp vụ trở thành một yêu cầu quan trọng. Kiến trúc phần mềm hướng dịch vụ, đặc biệt là kiến trúc microservices kết hợp với tư duy thiết kế theo miền nghiệp vụ, mang lại nhiều lợi ích trong việc phân tách trách nhiệm, tổ chức hệ thống và triển khai độc lập các thành phần.

Bài tiểu luận này tập trung tìm hiểu cơ sở lý thuyết về kiến trúc nguyên khối, kiến trúc microservices và Domain Driven Design; từ đó vận dụng vào việc phân tích, thiết kế hệ thống thương mại điện tử. Nội dung bài làm cũng đề cập đến vai trò của dịch vụ trí tuệ nhân tạo trong việc hỗ trợ gợi ý sản phẩm và tư vấn mua sắm, góp phần nâng cao trải nghiệm người dùng.

Do giới hạn về thời gian và phạm vi nghiên cứu, bài tiểu luận khó tránh khỏi những thiếu sót. Sinh viên thực hiện rất mong nhận được sự góp ý của giảng viên để hoàn thiện hơn nội dung nghiên cứu và rút kinh nghiệm cho các bài học, đồ án tiếp theo.
\end{document}